\begin{abstract}
流匹配已成为视觉运动机器人策略的强大生成式框架，具有训练稳定和高质量多模态动作生成等优势。然而，现有的流匹配策略在视觉编码阶段采用全局平均池化，将所有空间结构压缩为单一向量，丢失了精密任务所需的细粒度空间信息——策略知道它\emph{看到了什么}，却不知道交互\emph{发生在哪里}。

我们提出\textbf{\fasafull{}（\fasa{}）}，一个通过帧间特征时间对比自监督发现任务相关空间区域的即插即用模块。\fasa{}的核心洞察是：连续帧之间变化最显著的区域对应于当前活跃的交互区域。基于这一洞察，\fasa{}将空间特征分解为动态流（交互区域）和静态流（稳定上下文），编码变化质心和运动方向，并提取空间关键点坐标——提供比全局池化丰富$2\times$的空间信息，无需任何外部教师模型。结合特征差分编码的"语义速度"，\fasa{}为流匹配策略提供了丰富的空间感知能力。

\fasa{}不需要人工空间标注，无需外部教师模型，并能无缝集成到现有流匹配架构中。我们将其集成到VITA流匹配策略中，在一个具有挑战性的6自由度机器人S型焊缝跟踪任务上进行评估，实验表明\fasa{}显著提升了跟踪精度。
\end{abstract}
